\chapter{Fazit und Ausblick}
\autor{Sophie Kleppe}

\acrshort{bci}-Systeme können in Zukunft zu einer wegweisenden Technologie werden, die es Menschen mit Lähmungen ermöglicht, \glspl{neuroprothese} anzusteuern. Da der Einsatz im Alltag und als Experimentierstation im DLR\_School\_Lab allerdings nur mit einem portablen \acrshort{eeg}-System sinnvoll ist, war das Ziel dieser Arbeit zu prüfen, ob das OpenBCI-System auch in einer alltagsnahen Umgebung eine mit dem BioSemi-System vergleichbare Datenqualität liefert und damit als Grundlage für die Steuerung eines Roboterarms dienen kann. Um die Fragestellung zu beantworten, wurden zehn Versuchspersonen unter drei verschiedenen Versuchsbedingungen gemessen: Versuchsbedingung s1 (OpenBCI, Kabinentür offen, ohne Kinnstütze), s2 (OpenBCI, Kabinentür geschlossen, mit Kinnstütze) sowie s3 (BioSemi als Referenzsystem). Innerhalb der Versuchsbedingungen wurden zwei Paradigmen untersucht: das Alpha-Paradigma, bei dem die Augen abwechselnd geöffnet und geschlossen wurden, und das P1-Paradigma, bei dem Spielkarten groß und kontrastreich oder klein und kontrastreduziert gezeigt wurden.

Bezüglich des Alpha-Paradigmas zeigte sich, dass der \gls{berger-effekt} mit dem OpenBCI-System in allen Versuchsbedingungen deutlich messbar war, auch bei offener Kabinentür. Zudem lag die Alpha-Suppression bei allen Versuchspersonen über dem Kriterium von 2, die Alpha-Power war bei geschlossenen Augen also mindestens doppelt so hoch wie bei geöffneten. Sie war okzipito-parietal am stärksten. Die Rangfolge der Versuchspersonen stimmte stark mit der des BioSemi-Systems überein, bei offener Kabinentür stärker als bei geschlossener. Zurückzuführen ist dies vermutlich auf die zunehmende Müdigkeit der Versuchspersonen während der Messung in der geschlossenen Kabine.

Hinsichtlich des P1-Paradigmas war die P1 mit dem OpenBCI-System erkennbar und hob sich im Mittel vom Rauschen ab. Sie wies allerdings eine kleinere Amplitude auf und war stärker verrauscht. Eine mögliche Ursache ist die ungenaue Rekonstruktion der Trigger. Der Unterschied zwischen großen und kleinen Reizen war nur beim BioSemi-System belastbar. Beim OpenBCI-System lag er im Bereich des Messfehlers und kann deshalb nur als Tendenz aufgefasst werden. In der geschlossenen Kabine war die P1-Datenqualität besser als bei offener Kabinentür. Ob die Augen geöffnet oder geschlossen waren, wurde mit dem OpenBCI-System in 2-s-Abschnitten in mindestens 75\,\% der Fälle richtig erkannt, ein Ergebnis, das deutlich über der \gls{zufallsgrenze} liegt.

Für die Beantwortung der Fragestellung ergibt sich daraus, dass das OpenBCI-System für große und robuste Signale auch unter alltagsnahen Bedingungen eine ähnliche Datenqualität wie das BioSemi-System liefert. Für kleine Unterschiede in den \acrshortpl{erp} gilt dies nur eingeschränkt. Da das Öffnen und Schließen der Augen ein einfaches Steuersignal ist, das jede Person willentlich und ohne Training erzeugen kann, und bereits in Abschnitten von zwei Sekunden zuverlässig erkannt wird, kann das OpenBCI-System eine Grundlage für die Steuerung eines Roboterarms sein. Auch der Einsatz als Experimentierstation im DLR\_School\_Lab ist grundsätzlich möglich, da die Datenqualität für dieses Signal auch ohne abgeschirmte Kabine ausreicht. Allerdings gibt es auch Einschränkungen. Der Kontakt der Trockenelektroden ist nicht immer zuverlässig, sodass häufig Elektroden ausfielen und einzelne Messungen nicht auswertbar waren. Für das DLR\_School\_Lab bedeutet das, dass bei einer Station mit dem Ausfall einzelner Elektroden gerechnet werden muss. Zudem sollte die Anwendung anregend gestaltet sein, da Müdigkeit die Alpha-Aktivität beeinflusst.

Als nächster Schritt kann der Roboterarm mit dem OpenBCI-System in Echtzeit angesteuert werden. Dafür müsste die Alpha-Power während der Messung fortlaufend in kurzen Abschnitten berechnet und mit einer Schwelle verglichen werden. Wird die Schwelle überschritten, also werden die Augen geschlossen, könnte eine festgelegte Bewegung des Roboterarms ausgelöst werden. Da die Alpha-Suppression zwischen den Versuchspersonen stark schwankte, sollte die Schwelle vorher in einer kurzen Messung mit geöffneten und geschlossenen Augen für jede Person einzeln festgelegt werden. In einer Folgestudie müsste geprüft werden, ob diese Steuerung auch während der Messung zuverlässig und schnell genug funktioniert.

Außerdem ist ein reines Ein-Aus-Signal, das durch das Schließen und Öffnen der Augen erzeugt wird, zwar ein erster Schritt in Richtung einer \acrshort{bci}-gesteuerten Bewegung, reicht aber nicht für eine gezielte Greifbewegung aus. Hierzu müsste eine Folgestudie untersuchen, ob sich auch andere Signale wie die sensomotorischen Rhythmen (vgl. Abschnitt~\ref{sec:bcibaender}) für die Steuerung eignen. Diese Arbeit zeigt damit, dass bereits ein kostengünstiges, portables \acrshort{eeg}-System die Grundlage für eine einfache \acrshort{bci}-Steuerung bieten kann, im Bildungsbereich und langfristig vielleicht auch für Menschen, die ihre Gliedmaßen nicht mehr selbst bewegen können.
